% =====================================================================
%  Capítulo 1 — contenido
%
%  Este archivo contiene ÚNICAMENTE el cuerpo del capítulo (sin preámbulo,
%  sin \begin{document}). Lo consumen dos documentos:
%     - cap01.tex   -> entrega individual (portada + índices + referencias)
%     - libro.tex   -> libro completo unificado
%  Por eso nunca debe llevar portada, índices ni \printbibliography.
%
%  Convenciones:
%   * Figuras: guárdalas en capitulos/cap01/figuras/ con el prefijo cap01-
%     y cítalas sólo por su nombre: \includegraphics{cap01-esquema.png}
%   * Etiquetas: \label{fig:cap01:...}, \label{tab:cap01:...},
%                \label{sec:cap01:...}, \label{lst:cap01:...}
%   * Referencias: \cite{clave} usando las claves de bib/referencias.bib
% =====================================================================

\chapter{\TituloCapUno}
\label{cap:01}

\section{Introducción}
\label{sec:cap01:introduccion}

Escribe aquí la introducción del capítulo. Las citas se insertan con el comando
\verb|\cite| y se renderizan en formato IEEE, por ejemplo \cite{grama:intro-parallel:2003}.
También puedes citar varias fuentes a la vez
\cite{pacheco:intro-parallel-prog:2011, foster:disenando:1995}.

\section{Desarrollo}
\label{sec:cap01:desarrollo}

Texto del desarrollo. La taxonomía clásica de arquitecturas fue propuesta por
Flynn \cite{flynn:taxonomia:1972}, mientras que los límites teóricos del
paralelismo se discuten en \cite{amdahl:ley:1967} y \cite{gustafson:reevaluando:1988}.

% --------------------- Ejemplo de figura ---------------------
% Sustituye el \fbox por: \includegraphics[width=0.6\textwidth]{cap01-ejemplo.png}
\begin{figure}[H]
    \centering
    \fbox{\parbox[c][4cm][c]{0.6\textwidth}{\centering\itshape
        Marcador de posición: reemplázalo por\\
        \texttt{\textbackslash includegraphics\{cap01-ejemplo.png\}}}}
    \caption{Ejemplo de figura del capítulo 1. Todo lo que lleve \texttt{\textbackslash caption}
             aparece automáticamente en el Índice de figuras.}
    \label{fig:cap01:ejemplo}
\end{figure}

Como se aprecia en la Figura~\ref{fig:cap01:ejemplo}, las referencias cruzadas se
hacen con \verb|\ref| sobre la etiqueta declarada en el entorno.

% --------------------- Ejemplo de tabla ---------------------
\begin{table}[H]
    \centering
    \caption{Ejemplo de tabla del capítulo 1. Todo lo que lleve \texttt{\textbackslash caption}
             aparece automáticamente en el Índice de tablas.}
    \label{tab:cap01:ejemplo}
    \begin{tabular}{lrr}
        \toprule
        \textbf{Hilos} & \textbf{Tiempo (s)} & \textbf{Aceleración} \\
        \midrule
        1 & 12.40 & 1.00 \\
        2 &  6.85 & 1.81 \\
        4 &  3.72 & 3.33 \\
        8 &  2.15 & 5.77 \\
        \bottomrule
    \end{tabular}
\end{table}

La Tabla~\ref{tab:cap01:ejemplo} muestra el formato sugerido para reportar mediciones.

% --------------------- Ejemplo de código ---------------------
\begin{lstlisting}[style=C, caption={Ejemplo mínimo de una región paralela con OpenMP.},
                   label={lst:cap01:ejemplo}]
#include <stdio.h>
#include <omp.h>

int main(void) {
    #pragma omp parallel
    {
        printf("Hola desde el hilo %d de %d\n",
               omp_get_thread_num(), omp_get_num_threads());
    }
    return 0;
}
\end{lstlisting}

\begin{nota}
    Las cajas \texttt{nota} y \texttt{definicion} están definidas en
    \texttt{config/preambulo.tex} y sirven para resaltar ideas clave.
\end{nota}

\section{Conclusiones del capítulo}
\label{sec:cap01:conclusiones}

Cierra aquí las ideas principales del capítulo.
